
\chapter[双曲型演化方程与 Egorov 定理]{双曲型演化方程与\\Egorov 定理}\label{ch:07}
\index{shuangquxing yanhua fangcheng@双曲型演化方程 (hyperbolic evolution equations)}%
\index{Egorov dingli@Egorov 定理 (Egorov's theorem)}%

作为迄今所发展的伪微分算子机制的一项简洁应用, 我们证明双曲型演化方程解的存在性与唯一性 (第~\ref{sec:7.1}~节); 例子包括输运方程与波动方程. 我们还证明: 初始数据的奇性 (在波前集的意义下) 沿支配该演化方程的算子主符号的 Hamilton 流传播 (第~\ref{sec:7.2}~节). 我们只在 $\R^n$ 上陈述结果, 而把迁移到紧流形上所需的纯记号上的改动留给读者.

\section{存在性与唯一性}\label{sec:7.1}
\index{duichen shuangqu xitong@对称双曲系统 (symmetric hyperbolic systems)}%
\index{chuzhi wenti@初值问题 (initial value problems)}%

设 $a(t, x, \xi) \in \Smooth(\R_t; S^1(\R^n; \R^n)^{K \times K})$ 是一阶符号构成的 $K \times K$ 矩阵, 光滑依赖于 $t$; 我们简记为 $a \in S^1$. 令 $A(t) \defeq \Op_L(a(t))$. 考虑初值问题
\begin{equation}\label{eq:7.1}
    \left\{
    \begin{aligned}
        (D_t + A(t)) u &= g(t, x) \in \mathcal{C}^0\left( \R; H^s(\R^n; \CC^K) \right), \\
        u(0, x) &= u_0(x) \in H^s(\R^n; \CC^K).
    \end{aligned}
    \right.
\end{equation}
我们假定~\eqref{eq:7.1} 是\textbf{对称双曲的} (symmetric hyperbolic), 意即
\begin{equation}\label{eq:7.2}
    a(t, x, \xi) - a(t, x, \xi)^* \in S^0,
\end{equation}
或等价地 $A(t) - A(t)^* \in \Psi^0$.

\begin{theorem}[一阶对称双曲系统解的存在性与唯一性]\label{thm:7.1}
初值问题~\eqref{eq:7.1} 有唯一解
\begin{equation}\label{eq:7.3}
    u \in \mathcal{C}^0\left( \R; H^s(\R^n; \CC^K) \right) \cap \mathcal{C}^1\left( \R; H^{s-1}(\R^n; \CC^K) \right).
\end{equation}
\end{theorem}

\begin{proof}
我们在记号中略去``丛'' $\CC^K$. 我们将通过正则化方程的解 $u_\epsilon$ 的极限来得到 $u$:
\begin{equation}\label{eq:7.4}
    \left\{
    \begin{aligned}
        D_t u_\epsilon &= -J_\epsilon A J_\epsilon u_\epsilon + g, \\
        u_\epsilon(0) &= u_0,
    \end{aligned}
    \right.
\end{equation}
其中我们使用如~\eqref{eq:4.124} 中的 Friedrichs 磨光子 $\phi(\epsilon D) = J_\epsilon = J_\epsilon^*$, 于是当 $\epsilon > 0$ 时 $J_\epsilon \in \Psi^{-\infty}(\R^n)$, 且 $J_\epsilon \in \Psi^0(\R^n)$ 对 $\epsilon \in (0, 1]$ 一致有界, 并且对一切 $m' > 0$ 有 $J_\epsilon \to I$ 于 $\Psi^{m'}(\R^n)$ 中. 对 $\epsilon > 0$, $J_\epsilon A J_\epsilon$ 是 Banach 空间 $H^s(\R^n)$ 上关于 $t$ 连续的有界算子族, 因此~\eqref{eq:7.4} 在 $u_\epsilon \in \mathcal{C}^1(\R; H^s(\R^n))$ 中的可解性由标准的常微分方程理论得到.

我们需要建立关于 $u_\epsilon$ 的、当 $\epsilon \to 0$ 时一致的估计. 令 $\Lambda^s = \langle D_x \rangle^s$. 则
\begin{equation}\label{eq:7.5}
    \begin{aligned}
        \frac{1}{2} \frac{\mathrm{d}}{\mathrm{d} t} \| u_\epsilon(t) \|_{H^s}^2
        &= -\operatorname{Re} \langle \Lambda^s J_\epsilon i A J_\epsilon u_\epsilon, \Lambda^s u_\epsilon \rangle
         + \operatorname{Re} \langle \Lambda^s i g, \Lambda^s u_\epsilon \rangle \\
        &= -\operatorname{Re} \langle i A \Lambda^s J_\epsilon u_\epsilon, \Lambda^s J_\epsilon u_\epsilon \rangle
         - \operatorname{Re} \langle [\Lambda^s, i A] J_\epsilon u_\epsilon, \Lambda^s J_\epsilon u_\epsilon \rangle \\
        &\quad + \operatorname{Re} \langle \Lambda^s i g, \Lambda^s u_\epsilon \rangle.
    \end{aligned}
\end{equation}
由于 $B(t) = A(t) - A(t)^* \in \Psi^0$, 第一项被下式界住:
\begin{equation}\label{eq:7.6}
    \left| \langle B(t) \Lambda^s J_\epsilon u_\epsilon, \Lambda^s J_\epsilon u_\epsilon \rangle \right|
    \le C(t) \| J_\epsilon u_\epsilon \|_{H^s}^2
    \le \tilde{C}(t) \| u_\epsilon \|_{H^s}^2.
\end{equation}
由于 $[\Lambda^s, A] \in \Psi^s$, \eqref{eq:7.5} 中的第二项同样被 $C(t) \| u_\epsilon \|_{H^s}^2$ 界住. 对~\eqref{eq:7.5} 中的第三项应用 Cauchy--Schwarz 不等式, 我们得到 (取另一个常数 $C(t)$)
\begin{equation}\label{eq:7.7}
    \frac{\mathrm{d}}{\mathrm{d} t} \| u_\epsilon(t) \|_{H^s}^2
    \le C(t) \| u_\epsilon(t) \|_{H^s}^2 + \| g(t) \|_{H^s}^2.
\end{equation}
这蕴含与 $\epsilon$ 无关的估计
\begin{equation}\label{eq:7.8}
    \| u_\epsilon(t) \|_{H^s}^2
    \le \tilde{C}(t) \left( \| u_0 \|_{H^s}^2 + \| g \|_{\mathcal{C}^0([0, t]; H^s(\R^n))}^2 \right).
\end{equation}
固定 $T > 0$ 并令 $I \defeq [-T, T]$. 则~\eqref{eq:7.8} 蕴含 $u_\epsilon$ 在 $\mathcal{C}^0(I; H^s(\R^n))$ 中一致有界. 方程~\eqref{eq:7.4} 以及 $J_\epsilon A J_\epsilon \in \mathcal{L}(H^s(\R^n), H^{s-1}(\R^n))$ 的一致有界性蕴含
\begin{equation}\label{eq:7.9}
    u_\epsilon \in \mathcal{C}^0(I; H^s(\R^n)) \cap \mathcal{C}^1(I; H^{s-1}(\R^n)) \text{ 一致有界.}
\end{equation}
我们利用下面的引理~\ref{lem:7.2} 抽取 $u_\epsilon$ 的弱子列极限 $u \in \mathcal{C}^0(I; H^{s-1}(\R^n))$. 对该极限, 有 $u_\epsilon(0) = u_0 \rightharpoonup u(0)$ 于 $H^{s-1}(\R^n)$ 中, 故 $u(0) = u_0$. 此外, $u$ 在如下弱意义下满足偏微分方程~\eqref{eq:7.1}: 对一切 $\phi \in \SmoothComp(I^\circ; H^{-s+2}(\R^n))$,
\begin{equation}\label{eq:7.10}
    \begin{aligned}
        \langle D_t u, \phi \rangle
        &= \int_I \langle D_t u(t), \phi(t) \rangle_{L^2(\R^n)} \dd t
         = \langle u, -D_t \phi \rangle = \lim_{\epsilon \to 0} \langle u_\epsilon, -D_t \phi \rangle \\
        &= \lim_{\epsilon \to 0} \langle D_t u_\epsilon, \phi \rangle
         = \lim_{\epsilon \to 0} \langle J_\epsilon A J_\epsilon u_\epsilon, \phi \rangle \\
        &= \lim_{\epsilon \to 0} \langle u_\epsilon, J_\epsilon A^* J_\epsilon \phi \rangle
         = \langle u, A^* \phi \rangle = \langle A u, \phi \rangle,
    \end{aligned}
\end{equation}
其中在过渡到最后一行的过程中我们用到了 $J_\epsilon A^* J_\epsilon \phi \to A^* \phi$ 于 $\SmoothComp(I^\circ; H^{-s+1}(\R^n))$ 中 (习题~\ref{ex:4.8}). 于是 $u$ 是~\eqref{eq:7.1} 的弱解, 从而 $u \in \mathcal{C}^1(I; H^{s-2}(\R^n))$. $u$ 在空间 $\mathcal{C}^0(\R; H^\sigma(\R^n))$ 中 (对任意 $\sigma \in \R$) 的唯一性可由下述观察得到: 当 $u_0, g = 0$ 时, 导致~\eqref{eq:7.8} 的计算给出
\begin{equation*}
    \frac{\mathrm{d}}{\mathrm{d} t} \| u(t) \|_{H^{\sigma-1}}^2 \le C(t) \| u(t) \|_{H^{\sigma-1}}^2;
\end{equation*}
初值条件 $\| u(0) \|_{H^{\sigma-1}} = \| u_0 \|_{H^{\sigma-1}} = 0$ 于是蕴含对一切 $t$ 有 $u(t) = 0$.

为证明 $u$ 具有正确的正则性, 我们在上述拓扑中用 $u_{0, j} \in H^{s+1}(\R^n)$, $g_j \in \mathcal{C}^0(\R; H^{s+1}(\R^n))$, $j \in \N$, 逼近 $u_0 \in H^s(\R^n)$, $g \in \mathcal{C}^0(\R; H^s(\R^n))$. 此前的论证于是给出~\eqref{eq:7.1} 的解 $u_j \in \mathcal{C}^0(I; H^s(\R^n)) \cap \mathcal{C}^1(I; H^{s-1}(\R^n))$, 其中以 $u_{0, j}, g_j$ 代替 $u_0, g$. 此外,
\begin{equation}\label{eq:7.11}
    v_{jk} \defeq u_j - u_k
\end{equation}
解~\eqref{eq:7.1}, 其初值为 $u_{0, j} - u_{0, k}$ (在 $H^s(\R^n)$ 中 $\to 0$), 源项为 $g_j - g_k$ (在 $\mathcal{C}^0(\R; H^s(\R^n))$ 中 $\to 0$). 与~\eqref{eq:7.8} 类似的估计蕴含当 $j, k \to \infty$ 时 $v_{jk} \to 0$ 于 $\mathcal{C}^0(I; H^s(\R^n))$ 中. 因此 $\{u_j\}$ 是 $\mathcal{C}^0(I; H^s(\R^n))$ 中的 Cauchy 列, 从而其极限 $u$ 满足~\eqref{eq:7.3}.
\end{proof}

为完成证明, 我们需要:

\begin{lemma}[弱子列极限]\label{lem:7.2}
设 $(H, \langle \cdot, \cdot \rangle_H)$ 是 Hilbert 空间, 范数为 $\| \cdot \|_H$. 设 $a, b \in \R$, $a < b$, 并以 $\mathcal{C}^{0,1}([a, b]; H)$ 表示 Lipschitz 连续函数 $u \colon [a, b] \to H$ 所成之空间, 赋以范数
\begin{equation}\label{eq:7.12}
    \| u \|_{\mathcal{C}^{0,1}} \defeq \sup_{t \in [a, b]} \| u(t) \|_H + \sup_{s \neq t} \frac{\| u(t) - u(s) \|_H}{| t - s |}.
\end{equation}
设 $(u_j)_{j \in \N}$ 是 $\mathcal{C}^{0,1}([a, b]; H)$ 中的有界序列. 则存在子列 $(u_{j_k})_{k \in \N}$ 收敛到某个 $u \in \mathcal{C}^{0,1}([a, b]; H)$, 收敛的意义如下:
\begin{enumerate}
    \item 对一切 $t \in [a, b]$, 有 $u_{j_k}(t) \rightharpoonup u(t)$.
    \item 对一切 $\phi \in \mathcal{C}^{0,1}([a, b]; H)$, 有
    \begin{equation*}
        \langle u_{j_k}, \phi \rangle \defeq \int_a^b \langle u_{j_k}(t), \phi(t) \rangle_H \dd t \to \langle u, \phi \rangle.
    \end{equation*}
\end{enumerate}
\end{lemma}

\begin{proof}
我们作平移与伸缩使得 $[a, b] = [0, 1]$. 固定可数稠密子集 $\{t_i\} \subset [0, 1]$. 由于 $u_j(t_1)$ 是 $H$ 中的有界序列, 它有弱收敛子列 $u_{j_{1, k}}(t_1)$, 其弱极限记为 $\alpha_1 \in H$. 序列 $u_{j_{1, k}}(t_2)$ 有弱收敛子列 $u_{j_{2, k}}(t_2)$, 其弱极限记为 $\alpha_2 \in H$; 依此类推. 对对角子列 $u_{j_k} \defeq u_{j_{k, k}}$, 于是对一切 $i$ 有 $u_{j_k}(t_i) \rightharpoonup \alpha_i$.

今令 $C \defeq \sup_{j \in \N} \| u_j \|_{\mathcal{C}^{0,1}}$. 鉴于
\begin{equation}\label{eq:7.13}
    \| \alpha_i - \alpha_\ell \|_H
    \le \liminf_{k \to \infty} \| u_{j_k}(t_i) - u_{j_k}(t_\ell) \|_H
    \le C | t_i - t_\ell |,
\end{equation}
函数 $u \colon t_i \mapsto \alpha_i$ 具有 Lipschitz 常数 $\le C$ 的连续延拓. 若 $t \in [0, 1]$ 任意, $\phi \in H$, $\epsilon > 0$, 则
\begin{equation}\label{eq:7.14}
    \left| \langle u(t), \phi \rangle_H - \langle u_{j_k}(t), \phi \rangle_H \right|
    \le \left| \langle u(t_i) - u_{j_k}(t_i), \phi \rangle_H \right| + 2C | t - t_i | \| \phi \|_H < \epsilon,
\end{equation}
只要选取 $t_i$ 使 $2C | t - t_i | \| \phi \|_H < \frac{\epsilon}{2}$, 然后取 $k$ 充分大使得同时有 $\left| \langle u(t_i) - u_{j_k}(t_i), \phi \rangle_H \right| < \frac{\epsilon}{2}$.

最后, 给定 $\phi \in \mathcal{C}^{0,1}([0, 1]; H)$, 令 $L$ 表示 $[a, b] \ni t \mapsto \langle u(t), \phi(t) \rangle_H$ 的 Lipschitz 常数. 给定 $\epsilon > 0$, 我们有
\begin{equation}\label{eq:7.15}
    \left| \langle u, \phi \rangle - \frac{1}{N} \sum_{m=0}^{N-1} \langle u(m/N), \phi(m/N) \rangle \right| < \epsilon,
\end{equation}
只要 $\frac{L}{N} < \epsilon$, 即 $N > L \epsilon^{-1}$; 以 $u_{j_k}$ 代替 $u$ 亦类似. 固定这样的 $N$, 当 $k$ 充分大时对一切 $m = 0, \ldots, N-1$ 有 $\left| \langle u(m/N), \phi(m/N) \rangle_H - \langle u_{j_k}(m/N), \phi(m/N) \rangle_H \right| < \epsilon$. 因此 $\left| \langle u, \phi \rangle - \langle u_{j_k}, \phi \rangle \right| < 3\epsilon$. 这就证明了 $\langle u_{j_k}, \phi \rangle \to \langle u, \phi \rangle$.
\end{proof}

我们给出三个例子.

\begin{example}[输运方程]\label{eg:7.3}
设 $v \in \mathcal{C}^\infty_b(\R^n; \R^n)$. 取 $K = 1$ 且 $a(t, x, \xi) = v(x) \cdot \xi$, 则方程~\eqref{eq:7.1} 就是输运方程 $(\partial_t + v(x) \cdot \partial_x)u(t, x) = \ii g(t, x)$. 它可用特征线法求解: 若由 $\gamma_x(0) = x$, $\gamma_x'(t) = v(\gamma_x(t))$ 定义 $\gamma_x(t)$, 则
\begin{equation*}
    u(t, \gamma_x(t)) = u_0(x) + \int_0^t \ii g(s, \gamma_x(s)) \dd s.
\end{equation*}
依赖于时间的 $v$ 可类似处理.
\end{example}

\begin{example}[$\R^n$ 上的 Klein--Gordon 方程]\label{eg:7.4}
考虑初值问题
\begin{equation}\label{eq:7.16}
    \left\{
    \begin{aligned}
        (-D_t^2 + \Delta + 1) u(t, x) &= h(t, x) &&\in \mathcal{C}^0(\R; H^{s-1}(\R^n)), \\
        u(0, x) &= u_0(x) &&\in H^s(\R^n), \\
        D_t u(0, x) &= u_1(x) &&\in H^{s-1}(\R^n)
    \end{aligned}
    \right.
\end{equation}
此即 Klein--Gordon 方程. 令 $\Lambda \defeq \langle D \rangle$. 则问题~\eqref{eq:7.16} 等价于
\begin{equation}\label{eq:7.17}
    \left\{
    \begin{aligned}
        (D_t + A(t)) U &= G, \\
        U(0) &= U_0,
    \end{aligned}
    \right.
\end{equation}
其中
\begin{equation}\label{eq:7.18}
    U(t) \defeq \binom{\Lambda u}{D_t u}, \quad
    A \defeq \begin{pmatrix} 0 & -\Lambda \\ -\Lambda & 0 \end{pmatrix}, \quad
    U_0 \defeq \binom{\Lambda u_0}{u_1}, \quad
    G \defeq \binom{0}{-h}.
\end{equation}
注意 $U_0 \in H^{s-1}(\R^n; \CC^2)$, $G \in \mathcal{C}^0(\R; H^{s-1}(\R^n; \CC^2))$. 定理~\ref{thm:7.1} 给出~\eqref{eq:7.17} 的解 $U = (U_1, U_2) \in \mathcal{C}^0 H^{s-1} \cap \mathcal{C}^1 H^{s-2}$. 令 $u \defeq \Lambda^{-1} U_1 \in \mathcal{C}^0 H^s \cap \mathcal{C}^1 H^{s-1}$. $U$ 的方程的第一个分量蕴含 $U_2 = D_t u$. 因此 $u \in \mathcal{C}^0 H^s \cap \mathcal{C}^1 H^{s-1} \cap \mathcal{C}^2 H^{s-2}$ 解~\eqref{eq:7.16}.
\end{example}

\begin{example}[依赖于时间的 Riemann 流形上的波动方程]\label{eg:7.5}
设 $X$ 是闭流形, $h \in \Smooth(\R_t; \Smooth(X; S^2 T^* X))$ 是 $X$ 上 Riemann 度量的光滑族. 波动方程的初值问题
\begin{equation}\label{eq:7.19}
    \left\{
    \begin{aligned}
        (-D_t^2 + \Delta_{g(t)}) u(t, x) &= h(t, x) &&\in \mathcal{C}^0(\R; H^{s-1}(X)), \\
        u(0, x) &= u_0(x) &&\in H^s(X), \\
        D_t u(0, x) &= u_1(x) &&\in H^{s-1}(X)
    \end{aligned}
    \right.
\end{equation}
同样可以化成~\eqref{eq:7.1} 的形式 (以 $X$ 代替 $\R^n$). 首先取 $\Lambda_{\frac{1}{2}}(t) \in \Psi^{\frac{1}{2}}_{\mathrm{cl}}(X)$, 其主符号为 $| \xi |^{\frac{1}{2}}_{g(t)_x^{-1}}$, 且光滑依赖于 $t \in \R$. 然后令
\begin{equation}\label{eq:7.20}
    \Lambda(t) \defeq I + \Lambda_{\frac{1}{2}}(t)^* \Lambda_{\frac{1}{2}}(t) \in \Smooth(\R_t; \Psi^1_{\mathrm{cl}}(X)).
\end{equation}
该算子的主符号为 $| \xi |_{g(t)_x^{-1}}$. 此外它可逆, 逆算子为 $\Lambda(t)^{-1} \in \Smooth(\R_t; \Psi^{-1}_{\mathrm{cl}}(X))$. (这是构造~\eqref{eq:5.152} 的一个实例.) 问题~\eqref{eq:7.19} 于是等价于
\begin{equation}\label{eq:7.21}
    \left\{
    \begin{aligned}
        (D_t + A(t)) U &= G, \\
        U(0) &= U_0,
    \end{aligned}
    \right.
\end{equation}
其中
\begin{equation}\label{eq:7.22}
    \begin{gathered}
        U(t) \defeq \binom{\Lambda(t) u}{D_t u}, \qquad
        A(t) \defeq \begin{pmatrix} -(D_t \Lambda(t)) \circ \Lambda(t)^{-1} & -\Lambda(t) \\ -\Delta_{g(t)} \circ \Lambda(t)^{-1} & 0 \end{pmatrix}, \\
        U_0 \defeq \binom{\Lambda(0) u_0}{u_1} \in H^{s-1}(X; \CC^2), \qquad
        G \defeq \binom{0}{-h} \in \mathcal{C}^0\left( \R; H^{s-1}(X; \CC^2) \right).
    \end{gathered}
\end{equation}
定理~\ref{thm:7.1} 给出~\eqref{eq:7.21} 的解 $U = (U_1, U_2) \in \mathcal{C}^0 H^{s-1} \cap \mathcal{C}^1 H^{s-2}$. 函数 $u \defeq \Lambda(t)^{-1} U_1 \in \mathcal{C}^0 H^s \cap \mathcal{C}^1 H^{s-1} \cap \mathcal{C}^2 H^{s-2}$ 于是解~\eqref{eq:7.19}. --- 对 $\R_t \times \R^n_x$ 上波动方程 $(-D_t^2 + \Delta)u(t, x) = h$ 的初值问题可以同样方式处理, 此时取 $\Lambda(t) = \langle D \rangle$. 我们在第~\ref{ch:09}~章讨论一般波动方程的可解性及其解的正则性.
\end{example}



\section{Egorov 定理; 奇性的传播}\label{sec:7.2}
\index{Hamilton liu@Hamilton 流 (Hamiltonian flow)}%
\index{qixing chuanbo@奇性传播 (propagation of singularities)}%

我们开始在标量情形 (即~\eqref{eq:7.1} 中 $K = 1$) 下研究双曲型演化方程解的微局部行为. 于是考虑
\begin{equation}\label{eq:7.23}
    (D_t + A(t)) u(t, x) = 0,
\end{equation}
其中我们假定 $A \in \Smooth(\R_t; \Psi^1_{\mathrm{cl}}(\R^n))$ 是经典的. 用 $a \in \Smooth(\R_t; S^1_{\mathrm{hom}}(\R^n; \R^n \setminus \{0\}))$ 记 $A$ 的齐次主符号. (注意与上一节相比, 此处 $a$ 的含义有这点细微的改变.) 对称双曲性的要求意即
\begin{equation}\label{eq:7.24}
    a(t, x, \xi) \text{ 取实值.}
\end{equation}
我们用
\begin{equation}\label{eq:7.25}
    S(t, s) \colon u(s) \mapsto u(t)
\end{equation}
记~\eqref{eq:7.23} 的解算子. 由定理~\ref{thm:7.1}, 对一切 $\sigma \in \R$ 有 $S(t, s) \in \mathcal{L}(H^\sigma(\R^n))$. 此外 $S(t, s)$ 可逆, 且 $S(t, s)^{-1} = S(s, t)$.

\begin{theorem}[Egorov 定理]\label{thm:7.6}
设 $P_0 = \Op(p_0) \in \Psi^m(\R^n)$, 并定义
\begin{equation}\label{eq:7.26}
    P(t) \defeq S(t, 0) \circ P_0 \circ S(0, t), \qquad t \in \R.
\end{equation}
则 $P(t)$ 是一个 $m$ 阶伪微分算子加上一个光滑化算子: 存在算子
\begin{equation*}
    R \in \bigcap_{N \in \R} \Smooth\left( \R; \mathcal{L}(H^{-N}(\R^n), H^N(\R^n)) \right)
\end{equation*}
使得 $\tilde{P}(t) \defeq P(t) - R(t) \in \Psi^m(\R^n)$, 且光滑依赖于 $t \in \R$. $\tilde{P}(t)$ 的主符号由下式给出:
\begin{equation}\label{eq:7.27}
    \sigma^m(\tilde{P}(t))\left( \mathcal{C}(t, 0)(x_0, \xi_0) \right) = \sigma^m(P_0)(x_0, \xi_0),
\end{equation}
其中 $\mathcal{C}(t_1, t_0)$ 是依赖于时间的 Hamilton 向量场 $H_{a(t, x, \xi)}$ 从 $t_0$ 到 $t_1$ 的流; 即 $\mathcal{C}(t_1, t_0)(x_0, \xi_0) = \gamma(t_1)$, 其中 $\gamma(t_0) = (x_0, \xi_0)$, $\gamma'(s) = H_{a(s)}|_{\gamma(s)}$.
\end{theorem}

在闭流形 $X$ 上, 光滑化算子是残余伪微分算子. 因此, 对 $X$ 上标量双曲型方程而言, 定理~\ref{thm:7.6} 的类似结论就是 $P(t) \in \Psi^m(X)$.

我们先证明一个与~\eqref{eq:7.27} 相关的技术性结果:

\begin{lemma}[符号的输运]\label{lem:7.7}
设 $a \in \Smooth(\R_t; S^1(\R^n_x; \R^n_\xi))$. 再设 $m \in \R$, $q_0 \in S^m(\R^n; \R^n)$ 以及 $e \in \Smooth(\R; S^m(\R^n; \R^n))$. 设 $q = q(t, x, \xi)$ 是
\begin{equation}\label{eq:7.28}
    \left\{
    \begin{aligned}
        (\partial_t + H_{a(t)}) q(t, x, \xi) &= e(t, x, \xi), \\
        q(0, x, \xi) &= q_0(x, \xi)
    \end{aligned}
    \right.
\end{equation}
的解. 则 $q \in \Smooth(\R; S^m(\R^n; \R^n))$.
\end{lemma}

\begin{proof}
由
\begin{equation}\label{eq:7.29}
    \gamma_{t, x, \xi}(0) = (x, \xi), \qquad \gamma_{t, x, \xi}'(s) = H_{a(t+s)}|_{\gamma_{t, x, \xi}(s)}
\end{equation}
定义 $\gamma_{t, x, \xi}(s) = (X_{t, x, \xi}(s), \Xi_{t, x, \xi}(s))$. 取定 $t, x, \xi$ 的值, 并简记 $X = X_{t, x, \xi}$, $\Xi = \Xi_{t, x, \xi}$. 则方程~\eqref{eq:7.29} 读作
\begin{equation}\label{eq:7.30}
    \begin{aligned}
        X(0) &= x, \qquad X'(s) = (\partial_\xi a)(t + s, X(s), \Xi(s)), \\
        \Xi(0) &= \xi, \qquad \Xi'(s) = -(\partial_x a)(t + s, X(s), \Xi(s)).
    \end{aligned}
\end{equation}
当 $t, s$ 停留于 $\R$ 的有界子集中时, 存在常数 $C$ 使得
\begin{equation}\label{eq:7.31}
    | X'(s) | \le C, \qquad \frac{| \Xi'(s) |}{\langle \Xi(s) \rangle} \le C.
\end{equation}
(这里照例 $\langle \Xi(s) \rangle = (1 + | \Xi(s) |^2)^{\frac{1}{2}}$.) 这蕴含对新的常数 $C$ 有界:
\begin{equation}\label{eq:7.32}
    | X(s) - X(0) |, \quad \frac{\langle \Xi(s) \rangle}{\langle \Xi(0) \rangle}, \quad \frac{\langle \Xi(0) \rangle}{\langle \Xi(s) \rangle} \le C.
\end{equation}
特别地, $\gamma_{t, x, \xi}(s)$ 对一切 $t, s \in \R$ 与 $x, \xi \in \R^n$ 有定义.

回到常微分方程~\eqref{eq:7.28}, 我们可以把其唯一解写成
\begin{equation}\label{eq:7.33}
    \begin{aligned}
        q(t, x, \xi) &= q_0(X_{t, x, \xi}(-t), \Xi_{t, x, \xi}(-t)) \\
        &\quad + \int_0^t e(s, X_{t, x, \xi}(-t + s), \Xi_{t, x, \xi}(-t + s)) \dd s.
    \end{aligned}
\end{equation}
这定义了 $(t, x, \xi) \in \R \times \R^n \times \R^n$ 的光滑函数. 对任意固定的 $T > 0$, 我们将证明 $q \in \Smooth((-T, T); S^m(\R^n; \R^n))$. 界~\eqref{eq:7.32} 连同 $q_0, e$ 的符号界蕴含
\begin{equation}\label{eq:7.34}
    | q(t, x, \xi) | \le C \langle \xi \rangle^m, \qquad t \in (-T, T).
\end{equation}
为控制 $q$ 的导数, 我们首先需要控制 $\gamma_{t, x, \xi}(s)$ 关于 $x, \xi$ 的导数. 对~\eqref{eq:7.30} 关于 $z = x_i, \xi_i$ 求导给出方程组
\begin{equation}\label{eq:7.35}
    \begin{aligned}
        (\partial_z X)' &= \sum_{j=1}^n \left( (\partial_\xi \partial_{x_j} a) \partial_z X_j + (\partial_\xi \partial_{\xi_j} a) \partial_z \Xi_j \right), \\
        (\partial_z \Xi)' &= -\sum_{j=1}^n \left( (\partial_x \partial_{x_j} a) \partial_z X_j + (\partial_x \partial_{\xi_j} a) \partial_z \Xi_j \right).
    \end{aligned}
\end{equation}
$a$ 的符号界与界~\eqref{eq:7.32} 蕴含
\begin{equation}\label{eq:7.36}
    \left| (\partial_z X)' \right|, \; \frac{\left| (\partial_z \Xi)' \right|}{\langle \Xi \rangle}
    \le C \left( | \partial_z X | + \frac{| \partial_z \Xi |}{\langle \Xi \rangle} \right),
\end{equation}
只要 $t, s$ 停留于 $\R$ 的有界子集中. 积分之后, 这蕴含
\begin{equation}\label{eq:7.37}
    | \partial_z X |, \; \frac{| \partial_z \Xi |}{\langle \Xi \rangle}
    \le C \left( | \partial_z X(0) | + \frac{| \partial_z \Xi(0) |}{\langle \Xi(0) \rangle} \right).
\end{equation}
因此
\begin{equation}\label{eq:7.38}
    | \partial_{x_i} X | \le C, \quad
    | \partial_{x_i} \Xi | \le C \langle \Xi(0) \rangle, \quad
    | \partial_{\xi_i} X | \le C \langle \Xi(0) \rangle^{-1}, \quad
    | \partial_{\xi_i} \Xi | \le C.
\end{equation}
对~\eqref{eq:7.33} 关于 $x_i, \xi_i$ 求导并利用~\eqref{eq:7.32} 与~\eqref{eq:7.38}, 我们得到
\begin{equation}\label{eq:7.39}
    \left| \partial_{x_i} q(t, x, \xi) \right| \le C \langle \xi \rangle^m, \qquad
    \left| \partial_{\xi_i} q(t, x, \xi) \right| \le C \langle \xi \rangle^{m-1}, \qquad t \in (-T, T).
\end{equation}
由于由~\eqref{eq:7.28} 有 $\partial_t q = e - H_{a(t)} q$, 故这又蕴含 $| \partial_t q(t, x, \xi) | \le C \langle \xi \rangle^m$. 高阶导数以同样方式处理; 我们把这些细节验证留给读者.
\end{proof}

\begin{proof}[定理~\ref{thm:7.6} 的证明]
对~\eqref{eq:7.26} 关于 $t$ 求导给出方程
\begin{equation}\label{eq:7.40}
    P'(t) + \ii [A(t), P(t)] = 0, \qquad P(0) = P_0.
\end{equation}
我们将首先构造该伪微分算子值常微分方程的一个近似解 $Q(t) = \Op(q(t))$, $q(t, x, \xi) \in S^m$, 即
\begin{equation}\label{eq:7.41}
    Q'(t) + \ii [A(t), Q(t)] = R_1(t), \qquad Q(0) = P_0, \qquad R_1 \in \Smooth(\R_t; \Psi^{-\infty}(\R^n));
\end{equation}
然后我们将证明 $Q$ 与 $P$ 相差一个光滑化算子. 我们作拟设
\begin{equation}\label{eq:7.42}
    q(t) \sim \sum_{k=0}^\infty q_k(t), \qquad q_k(t) \in S^{m-k}.
\end{equation}
取~\eqref{eq:7.41} 的主符号给出
\begin{equation}\label{eq:7.43}
    (\partial_t + H_{a(t)}) q_0(t, x, \xi) = 0, \qquad q_0(0, x, \xi) = p_0(x, \xi).
\end{equation}
于是 $q_0(t, \mathcal{C}(t)(x, \xi)) = p_0(x, \xi)$. 引理~\ref{lem:7.7} 给出 $q_0 \in S^m$. 以迭代方式推进, 我们取 $q_j \in S^{m-j}$, $j \ge 1$, 为输运方程 (其初值由 $P_0$ 确定)
\begin{equation}\label{eq:7.44}
    (\partial_t + H_{a(t)}) q_j(t, x, \xi) = e_j(t, x, \xi)
\end{equation}
的解, 其中 $e_j \in S^{m-j}$ 是仅涉及 $A$ 的全符号与 $q_0, \ldots, q_{j-1}$ 的表达式.

如此安排好了~\eqref{eq:7.41} 之后, 我们现在证明: 对一切 $N \in \R$, 差 $R(t) = P(t) - Q(t)$ 把任意 $\tilde{f} \in H^{-N}(\R^n)$ 映入 $H^N(\R^n)$. 取 $\tilde{f} = S(t, 0)f$, 这等价于
\begin{equation}\label{eq:7.45}
    v(t) \defeq S(t, 0)P_0 f, \qquad w(t) \defeq Q(t)S(t, 0)f \quad \Longrightarrow \quad v(t) - w(t) \in H^N(\R^n).
\end{equation}
注意 $v(t)$ 与 $w(t)$ 解方程
\begin{equation}\label{eq:7.46}
    \begin{aligned}
        (D_t + A(t)) v &= 0, && v(0) = P_0 f, \\
        (D_t + A(t)) w &= -\ii R_1 S(t, 0) f, && w(0) = P_0 f.
    \end{aligned}
\end{equation}
因此, 令 $g(t) \defeq \ii R_1 S(t, 0) f \in \bigcap_{N \in \R} \Smooth(\R_t; H^N(\R^n))$, 我们有
\begin{equation}\label{eq:7.47}
    (D_t + A(t))(v - w) = g, \qquad (v - w)(0) = 0.
\end{equation}
定理~\ref{thm:7.1} 给出对一切 $N$ 有 $v - w \in \mathcal{C}^0(\R; H^N(\R^n))$. 由于 $D_t(v - w) = g - A(t)(v - w)$, 这又可以通过自举 (bootstrap) 提升为: 对一切 $N \in \R$, $k \in \N_0$, 有 $v - w \in \mathcal{C}^k(\R; H^N(\R^n))$.
\end{proof}

下面我们说明如何用定理~\ref{thm:7.6} 追踪标量演化方程解的波前集.

\begin{theorem}[波前集的传播: 标量情形]\label{thm:7.8}
设 $A$ 如~\eqref{eq:7.23}--\eqref{eq:7.24} 所述. 设 $s \in \R$. 设 $u_0 \in H^{-N}(\R^n)$, $g \in \mathcal{C}^0(\R_t; H^s(\R^n))$, 并设 $u$ 表示
\begin{equation}\label{eq:7.48}
    \left\{
    \begin{aligned}
        (D_t + A(t)) u &= g, \\
        u(0) &= u_0
    \end{aligned}
    \right.
\end{equation}
的解. 设 $\mathcal{C}(t_1, t_0)$ 如定理~\ref{thm:7.6} 所述. 则
\begin{equation}\label{eq:7.49}
    \WF^s(u(t)) = \mathcal{C}(t, 0)\WF^s(u_0),
\end{equation}
更一般地, 对一切 $t_0, t_1 \in \R$ 有 $\WF^s(u(t_1)) = \mathcal{C}(t_1, t_0)\WF^s(u(t_0))$.
\end{theorem}

\begin{proof}
从 $u$ 中减去 $(D_t + A(t))v = g$, $v(0) = 0$ 的解 $v \in \mathcal{C}^0(\R; H^s(\R^n))$, 我们可以化归到 $g = 0$ 的情形.

此时只需证明~\eqref{eq:7.49} 中的包含关系 ``$\subseteq$'', 因为把时间方向反过来便同时给出 ``$\supseteq$''. 于是设 $\alpha \notin \WF^s(u_0)$. 取算子 $P_0 \in \Psi^0(\R^n)$, 它在 $\alpha$ 处椭圆, 且使 $P_0 u_0 \in H^s(\R^n)$. 则在定理~\ref{thm:7.6} 的记号下
\begin{equation}\label{eq:7.50}
    P_0 u_0 = S(0, t)P(t)S(t, 0)u_0 \in H^s(\R^n),
\end{equation}
于是 $P(t)u(t) \in H^s(\R^n)$, 从而 $\tilde{P}(t)u(t) \in H^s(\R^n)$. 但 $\tilde{P}(t)$ 在 $\mathcal{C}(t, 0)\alpha$ 处椭圆, 故 $\mathcal{C}(t, 0)\alpha \notin \WF^s(u(t))$.
\end{proof}

\begin{remark}[齐次主符号]\label{rem:7.9}
上述论证不加改动地适用于具有实值齐次主符号 $a \in S^1_{\mathrm{hom}}$ 的一般算子 $A \in \Psi^1$. (也就是说, 若 $\chi \in \SmoothComp(\R^n_\xi)$ 在 $0$ 附近等于 $1$, 则 $(1 - \chi(\xi))a(t, x, \xi)$ 是 $\sigma^1(A(t))$ 的一个代表元.)
\end{remark}

\begin{remark}[推广到加权 Sobolev 空间]\label{rem:7.10}
定理~\ref{thm:7.1}, 从而定理~\ref{thm:7.8}, 可以 (只作纯记号上的改动) 推广到初值与源项属于加权 Sobolev 空间的情形. 借助定理~\ref{thm:2.18}, 于是可以允许定理~\ref{thm:7.8} 中的 $u_0$ 是任意缓增分布 $u_0 \in \Schwartz'(\R^n)$.
\end{remark}

\begin{example}[简单输运方程]\label{eg:7.11}
对 $A = D_x \in \Psi^1(\R)$, 解算子为 $(S(t, 0)u)(x) = u(x - t)$, 且 $\mathcal{C}(t, 0)(x_0, \xi_0) = (x_0 + t, \xi_0)$. 此外 $P(t) = \Op(p(t))$, $p(t, x, \xi) = p_0(x - t, \xi)$.
\end{example}

\begin{example}[半 Klein--Gordon 方程]\label{eg:7.12}
考虑 $\R_t \times \R^n_x$ 上的半 Klein--Gordon 方程
\begin{equation}\label{eq:7.51}
    (D_t + \langle D_x \rangle) u = 0.
\end{equation}
(它来自 Klein--Gordon 算子的分解 $-D_t^2 + (\Delta + 1) = -(D_t - \langle D_x \rangle)(D_t + \langle D_x \rangle)$.) 此时 $a(x, \xi) = | \xi |$, 其 Hamilton 向量场为 $H_a = \frac{\xi}{|\xi|} \cdot \partial_x$. 于是 $u$ 的波前集沿直线传播, 其方向由动量变量 $\xi$ 确定. (方程~\eqref{eq:7.51} 有显式解 $u(t) = \ee^{-\ii t \langle D \rangle}u(0)$, 波前集命题可由之直接验证. 当然, 定理~\ref{thm:7.6} 与~\ref{thm:7.8} 真正的力量在于它们同样适用于变系数方程.)
\end{example}



由于像 $\R \times \R^n$ 上波动方程或 Klein--Gordon 方程这样简单的例子并不被定理~\ref{thm:7.8} 覆盖, 有些读者可能希望走得更远. 我们进而证明一个一般性结果, 它确实允许控制一大类对称双曲系统解的波前集. 其做法是借助微局部解耦 (或称对角化) 论证化归到标量情形. (我们将在第~\ref{ch:08}~章联合 $t$ 与 $x$ 以及相应的典范动量变量, 重新考察波动方程解的波前集的传播: 与下面这些非标量论证相比, 那将是一个更简单的标量理论.)

对本节余下部分, 我们令
\begin{equation}\label{eq:7.52}
    A \in \Smooth\left( \R_t; \Psi^1_{\mathrm{cl}}(\R^n)^{K \times K} \right), \qquad
    a \defeq \sigma^1(A) \in \Smooth\left( \R; S^1_{\mathrm{hom}}(\R^n; \R^n \setminus \{0\})^{K \times K} \right).
\end{equation}

\begin{definition}[严格双曲性]\label{def:7.13}
我们说系统 $(D_t + A(t))u = g$ (或算子 $A$) 是\textbf{严格双曲的} (strictly hyperbolic), 若 $a(t, x, \xi)$ 对一切 $t \in \R$, $x \in \R^n$, $\xi \in \R^n \setminus \{0\}$ 有 $K$ 个互不相同的实特征值.
\end{definition}

我们的目标是把一个严格双曲系统解耦为 $K$ 个标量方程, 误差为光滑化算子. 我们用如下形式计算来引出这一构造: 若 $S(t)$ 是 $0$ 阶可逆算子, 则 $(D_t + A(t))u = 0$ 等价于 $S(t)u$ 的一个方程, 后者可通过如下改写得到:
\begin{equation}\label{eq:7.53}
    \begin{aligned}
        S(D_t + A)u &= S(D_t + A)S^{-1}Su \\
        &= \left( D_t + SAS^{-1} + S(D_tS^{-1}) \right)Su.
    \end{aligned}
\end{equation}
为使算子 $SAS^{-1} + S(D_tS^{-1})$ 为对角形, 我们首先需要 $SAS^{-1} \in \Psi^1_{\mathrm{cl}}$ 是对角形; 这是关于 $S$ 的主符号的线性代数问题. 项 $S(D_tS^{-1}) \in \Psi^0_{\mathrm{cl}}$ 阶数更低. 完全的解耦需要为 $S$ 构造合适的低阶修正.

\begin{proposition}[严格双曲系统的微局部解耦]\label{prop:7.14}
设 $A$ 是严格双曲的. 用
\begin{equation}\label{eq:7.54}
    \lambda_1(t, x, \xi) < \cdots < \lambda_K(t, x, \xi), \qquad t \in \R, \; x \in \R^n, \; \xi \in \R^n \setminus \{0\}
\end{equation}
记 $a(t, x, \xi)$ 的特征值. 设 $t_0 \in \R$, $x_0 \in \R^n$, $\xi_0 \in \R^n \setminus \{0\}$.
\begin{enumerate}
    \item (存在性.) 存在 $(t_0, x_0, \xi_0)$ 的一个 (关于 $\xi$ 的) 锥形邻域 $\mathcal{N} \subset \R_t \times \R^n_x \times (\R^n_\xi \setminus \{0\})$ 以及算子 $S \in \Smooth(\R; \Psi^0_{\mathrm{cl}}(\R^n)^{K \times K})$, 它在 $\overline{\mathcal{N}}$ 上椭圆, 其微局部参数子记为 $S^- \in \Smooth(\R; \Psi^0_{\mathrm{cl}}(\R^n)^{K \times K})$, 使得对算子
    \begin{equation}\label{eq:7.55}
        \mathcal{D} = (\mathcal{D}_{kl})_{k, l = 1, \ldots, K} \defeq S A S^- + S(D_t S^-) \in \Smooth\left( \R; \Psi^1_{\mathrm{cl}}(\R^n)^{K \times K} \right)
    \end{equation}
    有
    \begin{equation}\label{eq:7.56}
        \sigma^1(\mathcal{D}_{kk}) = \lambda_k \text{ 于 } \overline{\mathcal{N}} \text{ 附近}, \; k = 1, \ldots, K; \qquad
        \WF'(\mathcal{D}_{kl}) \cap \overline{\mathcal{N}} = \emptyset \;\; \forall\, k \neq l.
    \end{equation}
    \item (唯一性.) 若 $S, S'$ 是两个使第 (1) 部分的结论对同一个 $\mathcal{N}$ 成立的算子, 则存在算子 $B_1, \ldots, B_K \in \Smooth(\R; \Psi^0_{\mathrm{cl}}(\R^n))$, 它们在 $\overline{\mathcal{N}}$ 上椭圆, 使得
    \begin{equation}\label{eq:7.57}
        \WF'(S'_{kl} - B_k S_{kl}) \cap \overline{\mathcal{N}} = \emptyset, \qquad k, l = 1, \ldots, K.
    \end{equation}
\end{enumerate}
\end{proposition}

\begin{example}[Klein--Gordon 方程]\label{eg:7.15}
考虑把 Klein--Gordon 方程化为一阶 $2 \times 2$ 系统的约化~\eqref{eq:7.17}--\eqref{eq:7.18}. 则
\begin{equation}\label{eq:7.58}
    S \defeq \begin{pmatrix} 1 & 1 \\ 1 & -1 \end{pmatrix}
    \quad \Longrightarrow \quad
    \mathcal{D} \defeq S A S^{-1} + S(D_t S^{-1}) = \begin{pmatrix} -\Lambda & 0 \\ 0 & \Lambda \end{pmatrix}.
\end{equation}
此外 $\lambda_1(t, x, \xi) = -| \xi |$, $\lambda_2(t, x, \xi) = | \xi |$.
\end{example}

\begin{proof}[命题~\ref{prop:7.14} 的证明]
\textbf{存在性.} \emph{第 1 步: 对角化.} 由于 $a(t, x, \xi)$ 的特征值的代数重数为 $1$, 它们光滑依赖于 $t, x, \xi$. $a$ 关于 $\xi$ 的齐性蕴含 $\lambda_k(t, x, \xi) = | \xi | \lambda_k(t, x, \frac{\xi}{|\xi|})$. 因此 $\lambda_k$ 是 $(t_0, x_0, \xi_0)$ 的一个锥形邻域 $\mathcal{N}_0$ 中关于 $\xi$ 的 $1$ 阶齐次符号. 类似地, 到相应特征子空间的投影 $\pi_k(t, x, \xi)$ 是一个取值于 $\CC^{K \times K}$ 的 $0$ 阶齐次符号. 取 $v_k \in \CC^K$ 使 $\pi_k(t_0, x_0, \xi_0)v_k \neq 0$, 然后定义 $\CC^{K \times K}$ 值符号
\begin{equation}\label{eq:7.59}
    \tilde{s}_0(t, x, \xi) = (\tilde{s}_{0, kl}(t, x, \xi))_{k, l = 1, \ldots, K}, \qquad
    \tilde{s}_{0, kl}(t, x, \xi) \defeq (\pi_l(t, x, \xi)v_l)_k,
\end{equation}
它定义在 $\mathcal{N}_0$ 上, 是 $0$ 次齐次的. $\tilde{s}_0$ 的各列构成 $\CC^K$ 的一组由 $a(t, x, \xi)$ 的特征向量组成的基, 从而
\begin{equation}\label{eq:7.60}
    \tilde{s}_0^{-1} a \tilde{s}_0 = d_1 \defeq \operatorname{diag}(\lambda_1, \ldots, \lambda_K),
\end{equation}
其中 $\operatorname{diag}$ 表示以给定对角元为对角元的对角矩阵. 设 $\mathcal{N} \subset \mathcal{N}' \subset \mathcal{N}_0$ 是 $(t_0, x_0, \xi_0)$ 的锥形邻域, 且 $\overline{\mathcal{N}} \subset \mathcal{N}'$, $\overline{\mathcal{N}'} \subset \mathcal{N}_0$. 固定 $\chi \in S^0_{\mathrm{hom}}$, 使它在 $\overline{\mathcal{N}'}$ 附近等于 $1$, 在 $\mathcal{N}_0$ 的余集附近等于 $0$. 然后令
\begin{equation}\label{eq:7.61}
    s_0(t, x, \xi) \defeq \chi(t, x, \xi)\tilde{s}_0(t, x, \xi)^{-1}.
\end{equation}
取 $S_0 \in \Psi^0_{\mathrm{cl}}$ 使 $\sigma^0(S_0) = s_0$. 从今往后, 当 $A - B$ 等于一个符号、相应地一个算子的本质支集、相应地算子波前集与 $\overline{\mathcal{N}'}$ 不相交 (再加上 $X$ 中的一个元素) 时, 我们写作 ``$A \equiv B$'' (以及 ``$A \equiv B \bmod X$''). 令
\begin{equation}\label{eq:7.62}
    \mathcal{D}_1 = \operatorname{diag}(\Lambda_1, \ldots, \Lambda_K), \qquad
    \Lambda_k \in \Psi^1_{\mathrm{cl}}, \; \sigma^1(\Lambda_k) = \chi \lambda_k, \; k = 1, \ldots, K.
\end{equation}
由 $s_0$ 的构造, 我们于是有
\begin{equation}\label{eq:7.63}
    s_0 a s_0^{-1} = d_1 \text{ 于 } \overline{\mathcal{N}'} \text{ 上}, \qquad
    R_0 \defeq S_0 A S_0^- + S_0(D_t S_0^-) - \mathcal{D}_1 \equiv 0 \bmod \Psi^0_{\mathrm{cl}}.
\end{equation}
这里以及下文中, $S^-$ 表示算子 $S \in \Psi^1$ 在 $\overline{\mathcal{N}'}$ 上的微局部椭圆参数子.

\emph{第 2 步: 第一个低阶修正.} 我们的目标是证明: 若适当选取 $E_{-1} \in \Psi^{-1}_{\mathrm{cl}}$, 则
\begin{equation}\label{eq:7.64}
    S_1 A S_1^- + S_1(D_t S_1^-) - (\mathcal{D}_1 + \mathcal{D}_0) \equiv 0 \bmod \Psi^{-1}_{\mathrm{cl}}, \qquad S_1 \defeq (I + E_{-1})S_0,
\end{equation}
其中 $\mathcal{D}_0 \in \Psi^0_{\mathrm{cl}}$ 为某个对角算子. 为确定 $e_{-1} \defeq \sigma^{-1}(E_{-1}) \in S^{-1}_{\mathrm{hom}}$ 所需满足的方程, 首先注意 $S_1^- \equiv S_0^-(I - E_{-1}) \bmod \Psi^{-2}_{\mathrm{cl}}$. 于是取~\eqref{eq:7.64} 的 $0$ 阶主符号表明, 我们需要找 $e_{-1}$ 与 $d_0$ 使得
\begin{equation}\label{eq:7.65}
    r_0 + [e_{-1}, d_1] - d_0 = 0;
\end{equation}
这里 $r_0 \in S^0_{\mathrm{hom}}$ 是~\eqref{eq:7.63} 中误差项在 $\overline{\mathcal{N}'}$ 附近的主符号. (更确切地说, 取 $\chi \in S^0_{\mathrm{cl}}$ 在 $\overline{\mathcal{N}'}$ 附近等于 $1$ 且使 $\Op(\chi)R_0 \in \Psi^0_{\mathrm{cl}}$; 然后令 $r_0 \defeq \sigma^0(\Op(\chi)R_0)$.) 对固定的 $t, x, \xi$, 线性映射
\begin{equation}\label{eq:7.66}
    \CC^{K \times K} \ni e \mapsto [e, d_1(t, x, \xi)] = \left( (\lambda_l(t, x, \xi) - \lambda_k(t, x, \xi)) e_{kl} \right)_{k, l = 1, \ldots, K}
\end{equation}
的核等于对角矩阵所成之空间, 其值域等于对角元全为零的矩阵所成之空间. 因此我们可以在 $\overline{\mathcal{N}'}$ 附近如下解~\eqref{eq:7.65}: 先令 $d_0 = \operatorname{diag}((r_0)_{11}, \ldots, (r_0)_{KK}) \in S^0_{\mathrm{cl}}$; 然后令 $e_{-1} \in S^{-1}_{\mathrm{cl}}$ 为~\eqref{eq:7.65} 的满足 $(e_{-1})_{kk} = 0$ ($k = 1, \ldots, K$) 的唯一解. 取一个以 $\chi d_0$ 为主符号的对角算子 $\mathcal{D}_0$ 以及一个以 $\chi e_{-1}$ 为主符号的算子 $E_{-1}$; 这里 $\chi \in S^0_{\mathrm{cl}}$ 在 $\overline{\mathcal{N}'}$ 附近等于 $1$ 且支集位于 $\overline{\mathcal{N}'}$ 的一个小邻域中. 至此我们已安排好~\eqref{eq:7.64}.

\emph{第 3 步: 其余修正; 结论.} 以此方式推进, 对 $J = 2, 3, 4, \ldots$ 得到对角算子 $\mathcal{D}_{-J+1} \in \Psi^{-J+1}_{\mathrm{cl}}$ 与算子 $E_{-J} \in \Psi^{-J}_{\mathrm{cl}}$, 使得对 $S_J \defeq (I + E_{-J}) \circ \cdots \circ (I + E_{-1}) \circ S_0$ 有
\begin{equation}\label{eq:7.67}
    S_J A S_J^- + S_J(D_t S_J^-) - (\mathcal{D}_1 + \cdots + \mathcal{D}_{-J+1}) \equiv 0 \bmod \Psi^{-J}_{\mathrm{cl}}.
\end{equation}
然后选取 $E \in \Psi^{-1}_{\mathrm{cl}}$ 使得
\begin{equation}\label{eq:7.68}
    I + E - (I + E_{-J}) \circ \cdots \circ (I + E_{-1}) \in \Psi^{-J-1}_{\mathrm{cl}}
\end{equation}
对一切 $J \in \N$ 成立, 即 $I + E$ 是渐近乘积 $\cdots \circ (I + E_{-2}) \circ (I + E_{-1})$. 令 $S \defeq (I + E)S_0$, 并取对角算子 $\mathcal{D} \in \Psi^1_{\mathrm{cl}}$ 使 $\mathcal{D} \sim \mathcal{D}_1 + \mathcal{D}_0 + \cdots$. 则~\eqref{eq:7.67} 蕴含 $\WF'(SAS^- + S(D_t S^-) - \mathcal{D}) \cap \overline{\mathcal{N}} = \emptyset$. 这样我们便安排了~\eqref{eq:7.55}--\eqref{eq:7.56}.

\textbf{唯一性.} 变换
\begin{equation}\label{eq:7.69}
    \Phi \colon (A, S) \mapsto SAS^- + S(D_t S^-)
\end{equation}
具有性质 $\Phi(\Phi(A, S), T) \equiv \Phi(A, T \circ S)$. 令 $\mathcal{D} \defeq \Phi(A, S)$, $\mathcal{D}' \defeq \Phi(A, S')$, 则 $\Phi(\mathcal{D}, S'S^-) \equiv \mathcal{D}'$. 若 $\mathcal{D}$ 与 $\mathcal{D}'$ 都满足~\eqref{eq:7.56}, 我们需要证明: $S'S^-$ 在 $\overline{\mathcal{N}}$ 附近微局部地等于 $\operatorname{diag}(B_1, \ldots, B_K)$, 其中 $B_k \in \Psi^0_{\mathrm{cl}}$.

用新的记号, 我们因此必须证明: 给定一个严格双曲算子 $A$, 它以 $A_{kl}$ 代替 $\mathcal{D}_{kl}$ 且在 $\overline{\mathcal{N}'}$ 上 ($\mathcal{N}' \supset \overline{\mathcal{N}}$) 满足~\eqref{eq:7.56}, 并给定 $S \in \Psi^0_{\mathrm{cl}}$ 使 $\mathcal{D} = \Phi(A, S)$ 同样满足~\eqref{eq:7.56}, 则算子 $S$ 在 $\overline{\mathcal{N}}$ 附近微局部地对角.

首先注意 $A$ 与 $\mathcal{D}$ 在 $\overline{\mathcal{N}}$ 附近有相同的主符号 $\operatorname{diag}(\lambda_1, \ldots, \lambda_K)$. 由于 $A$ 的特征值 $\lambda_k$ 两两互异, 这迫使 $\sigma^0(S)$ 为对角矩阵. 设 $S_0 \in \Psi^0_{\mathrm{cl}}$ 是以 $\sigma^0(S)$ 为主符号的对角算子. 记 $S_0^- S =: I + \tilde{S}_{-1}$; 则 $\tilde{S}_{-1} \equiv 0 \bmod \Psi^{-1}_{\mathrm{cl}}$. 因此同样
\begin{equation}\label{eq:7.70}
    \begin{aligned}
        \Phi(\mathcal{D}, S_0^-) &\equiv \Phi(A, S_0^- S) \\
        &\equiv \Phi(A, I + \tilde{S}_{-1}) \\
        &\equiv (I + \tilde{S}_{-1})A(I + \tilde{S}_{-1})^- + (I + \tilde{S}_{-1})D_t\left( (I + \tilde{S}_{-1})^- \right) \\
        &\equiv A - [A, \tilde{S}_{-1}] \bmod \Psi^{-1}_{\mathrm{cl}}
    \end{aligned}
\end{equation}
是微局部对角的, 从而 $[A, \tilde{S}_{-1}]$ 也是. 回顾映射~\eqref{eq:7.66} 的性质, 这又迫使 $\sigma^{-1}(\tilde{S}_{-1})$ 为对角形. 设 $S_{-1} \in \Psi^{-1}_{\mathrm{cl}}$ 是以 $\sigma^{-1}(\tilde{S}_{-1})$ 为主符号的对角算子.

由于 $(I - S_{-1}) \circ (I + \tilde{S}_{-1}) =: I + \tilde{S}_{-2}$, 其中 $\tilde{S}_{-2} \equiv 0 \bmod \Psi^{-2}_{\mathrm{cl}}$, 我们现在断定
\begin{equation}\label{eq:7.71}
    \Phi(\mathcal{D}, (I - S_{-1}) \circ S_0^-) \equiv \Phi(A, I + \tilde{S}_{-2}) \equiv A - [A, \tilde{S}_{-2}] \bmod \Psi^{-2}_{\mathrm{cl}}
\end{equation}
是微局部对角的, 从而 $\sigma^{-2}(\tilde{S}_{-2})$ 也是; 依此类推. 以这种方式, 我们确定对角算子 $S_{-j} \in \Psi^{-j}_{\mathrm{cl}}$, $j = 2, 3, \ldots$, 使得
\begin{equation}\label{eq:7.72}
    (I - S_{-J+1}) \circ \cdots \circ (I - S_{-1}) \circ S_0^- \circ S =: I + \tilde{S}_{-J}, \qquad \tilde{S}_{-J} \equiv 0 \bmod \Psi^{-J}_{\mathrm{cl}}.
\end{equation}
取 $E \in \Psi^{-1}_{\mathrm{cl}}$ 使对一切 $J \in \N$ 有 $I + E - (I - S_{-J+1}) \circ \cdots \circ (I - S_{-1}) \in \Psi^{-J}_{\mathrm{cl}}$, 我们于是断定 $S \equiv S_0 \circ (I + E)^-$ 在 $\overline{\mathcal{N}}$ 附近微局部地对角.
\end{proof}

回到~\eqref{eq:7.53}, 在~\eqref{eq:7.55}--\eqref{eq:7.56} 之下, 现在的预期是: $S(t)u(t)$ 的每个分量 $(S(t)u(t))_k$, $k = 1, \ldots, K$, 在 $\overline{\mathcal{N}}$ 附近微局部地满足一个解耦的标量双曲型方程, 该方程由算子 $\mathcal{D}_{kk}$ 支配. 按~\eqref{eq:7.57} 的记号, 对一切 $u \in H^{-N}(\R^n; \CC^K)$ 在 $\overline{\mathcal{N}}$ 上有 $(S'(t)u)_k \equiv B_k(S(t)u)_k$, 其意义是两边之差的波前集与 $\overline{\mathcal{N}}$ 不相交. 由于算子 $B_k$ 在 $\overline{\mathcal{N}}$ 上椭圆, $(S(t)u(t))_k$ 在 $\overline{\mathcal{N}}$ 上的波前集因此不依赖于命题~\ref{prop:7.14}(1) 中算子 $S(t)$ 的选取. 这使我们可以无歧义地写
\begin{equation}\label{eq:7.73}
    \WF(u(t); k) \defeq \WF\left( (S(t)u(t))_k \right) \quad \text{于 } \overline{\mathcal{N}} \text{ 上}.
\end{equation}



\begin{theorem}[波前集的传播: 系统情形]\label{thm:7.16}
设 $u_0 \in H^{-N}(\R^n; \CC^K)$, $g \in \mathcal{C}^0(\R; H^\infty(\R^n; \CC^K))$, 并设 $u$ 是对称且严格双曲系统
\begin{equation}\label{eq:7.74}
    \left\{
    \begin{aligned}
        (D_t + A(t)) u &= g, \\
        u(0) &= u_0
    \end{aligned}
    \right.
\end{equation}
的解. 对 $k = 1, \ldots, K$, 设 $\mathcal{C}_k(t_1, t_0)$ 是依赖于时间的 Hamilton 向量场 $H_{\lambda_k(t)}$ 从 $t_0$ 到 $t_1$ 的流. 则
\begin{equation}\label{eq:7.75}
    \WF(u(t_1); k) = \mathcal{C}_k(t_1, t_0)\WF(u(t_0); k), \qquad k = 1, \ldots, K.
\end{equation}
特别地, 我们有
\begin{equation}\label{eq:7.76}
    \WF(u(t)) \subset \bigcup_{k=1}^K \mathcal{C}_k(t, 0)\WF(u_0).
\end{equation}
\end{theorem}

对称双曲性的假设可以去掉 (习题~\ref{ex:7.4}). \eqref{eq:7.75}--\eqref{eq:7.76} 的 Sobolev 精细化, 其陈述与证明我们留给有兴趣的读者.

\begin{proof}[定理~\ref{thm:7.16} 的证明]
如同定理~\ref{thm:7.8} 的证明, 只需考虑 $g = 0$ 的情形. 由定理~\ref{thm:7.1}, 方程~\eqref{eq:7.74} 有解 $u \in \mathcal{C}^0(\R; H^{-N}(\R^n)) \cap \mathcal{C}^1(\R; H^{-N-1}(\R^n))$.

\emph{化归为标量方程.} 在点 $(t_0, x_0, \xi_0) \in \R \times \R^n \times (\R^n \setminus \{0\})$ 处, 取 $S, S^-, \mathcal{N}$ 如命题~\ref{prop:7.14}. 我们的第一个目标是把 $u$ 的方程化归为关于 $S(t)u(t)$ 的某个分量的标量方程. 我们有 $S^- S = I - R$, 其中 $R \in \Psi^0_{\mathrm{cl}}$, $\WF'(R) \cap \overline{\mathcal{N}} = \emptyset$. 形式计算~\eqref{eq:7.53} 的严格版本于是读作
\begin{equation}\label{eq:7.77}
    \begin{aligned}
        0 = S(D_t + A)u &= S(D_t + A)(S^- S + R)u \\
        &= \left( D_t + S A S^- + S(D_t S^-) \right)Su + S(D_t + A)Ru.
    \end{aligned}
\end{equation}
令 $\mathcal{D} \defeq S A S^- + S(D_t S^-) \in \Psi^1_{\mathrm{cl}}$, $\tilde{h} \defeq -S(D_t + A)Ru$, 上式成为
\begin{equation}\label{eq:7.78}
    (D_t + \mathcal{D})(S(t)u(t)) = \tilde{h}.
\end{equation}
今设 $Q \in \Psi^0_{\mathrm{cl}}$ 是标量算子 (即它逐分量作用于 $\CC^K$ 值分布), 具有实主符号, 且 $\WF'(I - Q) \cap \overline{\mathcal{N}} = \emptyset$, 并使 $\WF'(Q)$ 含于 $\sigma^1(\mathcal{D}_{kk}) = \lambda_k$ 成立之集合中 (见~\eqref{eq:7.56}), 且对一切 $k \neq l$ 有 $\WF'(Q) \cap \WF'(\mathcal{D}_{kl}) = \emptyset$. 则
\begin{equation}\label{eq:7.79}
    (D_t + Q\mathcal{D})(Su) = h \defeq \tilde{h} - (I - Q)\mathcal{D}Su = -S(D_t + A)Ru - (I - Q)\mathcal{D}Su.
\end{equation}
最后, 我们考虑该方程的第 $k$ 个分量. 由于 $Q\mathcal{D}$ 的非对角项是残余算子, 我们得到标量对称双曲方程
\begin{equation}\label{eq:7.80}
    (D_t + Q\mathcal{D}_{kk})v(t) = h_{(k)}, \qquad v(t) \defeq (S(t)u(t))_k, \quad h_{(k)} - h_k \in \mathcal{C}^0(\R; H^\infty(\R^n)).
\end{equation}
由于 $h$ 从而 $h_{(k)}$ 的波前集与 $\overline{\mathcal{N}}$ 不相交 (因为 $\WF'(R)$ 与 $\WF'(I - Q)$ 都与 $\overline{\mathcal{N}}$ 不相交), 我们预期 $\WF(v(t))$ 沿 $\partial_t + H_{\lambda_k(t)}$ 的流传播, 至少在一个以 $(t_0, x_0, \xi_0)$ 为顶点的锥形邻域中如此.

\emph{波前集的传播.} 用 $S_k(t_1, t_0)$ 记齐次问题 $(D_t + Q\mathcal{D}_{kk})w(t) = 0$ 的解算子. 取 $\epsilon > 0$ 使下述成立: 对一切 $t \in [t_0 - \epsilon, t_0 + \epsilon]$ 与
\begin{equation}\label{eq:7.81}
    (x, \xi) \in \mathcal{N}_\epsilon(x_0, \xi_0) \defeq \left\{ (z, \zeta) \in \R^n \times (\R^n \setminus \{0\}) \colon | z - x_0 | < \epsilon, \; \left| \frac{\zeta}{|\zeta|} - \frac{\xi_0}{|\xi_0|} \right| < \epsilon \right\},
\end{equation}
对一切 $s \in [t_0 - \epsilon, t_0 + \epsilon]$ 有 $\mathcal{C}_k(t, s)(x, \xi) \in \mathcal{N}$. 今设 $(x_1, \xi_1) \in \mathcal{N}_\epsilon(x_0, \xi_0)$, $(x_1, \xi_1) \notin \WF(v(t_0))$. 取 $P_0 \in \Psi^0(\R^n)$ 在 $(x_1, \xi_1)$ 处椭圆, 使 $P_0 v(t_0) \in H^\infty(\R^n)$ 且 $\WF'(P_0) \subset \mathcal{N}_\epsilon(x_0, \xi_0)$. 定义
\begin{equation}\label{eq:7.82}
    P(t) \defeq S_k(t, t_0) \circ P_0 \circ S_k(t_0, t), \qquad t \in [t_0 - \epsilon, t_0 + \epsilon],
\end{equation}
这里用到定理~\ref{thm:7.6}. 由 Duhamel 公式, 我们可以把~\eqref{eq:7.80} 的解写成
\begin{equation}\label{eq:7.83}
    v(t) = S_k(t, t_0)v(t_0) + \int_{t_0}^t \ii S_k(t, s)h_{(k)}(s) \dd s.
\end{equation}
因此
\begin{equation}\label{eq:7.84}
    P(t)v(t) = S_k(t, t_0)\left( P_0 v(t_0) \right) + \int_{t_0}^t \ii S_k(t, s)\left( P(s)h_{(k)}(s) \right) \dd s.
\end{equation}
定理~\ref{thm:7.6} 的证明中对 $P(s)$ 的构造表明: $P(s)$, $s \in [t_0 - \epsilon, t_0 + \epsilon]$, 是一个光滑化算子与一个伪微分算子之和, 后者的算子波前集 $\WF'(P(s)) = \mathcal{C}_k(s, t_0)(\WF'(P_0))$ 是 $\mathcal{N}$ 的子集. (我们在此用到 $\mathcal{D}_{kk}$ 与 $Q\mathcal{D}_{kk}$ 的主符号在 $\mathcal{N}$ 上相等.) 因此~\eqref{eq:7.84} 的右端落在 $H^\infty(\R^n)$ 中, 从而 $\mathcal{C}_k(t, t_0)(x_1, \xi_1) \notin \WF(v(t))$, $t \in [t_0 - \epsilon, t_0 + \epsilon]$.

至此我们证明了: 当 $| t_1 - t_0 |$ 充分小时, \eqref{eq:7.75} 在 $\R^n \times (\R^n \setminus \{0\})$ 中每一点的某个锥形邻域内成立. 把这些局部结果拼接起来便给出对一切 $t_0, t_1$ 的~\eqref{eq:7.75}. 包含关系~\eqref{eq:7.76} 最终由
\begin{equation}\label{eq:7.85}
    \WF(u(t)) = \bigcup_{k=1}^K \WF(u(t); k), \qquad \WF(u(0); k) \subset \WF(u_0) \;\; (k = 1, \ldots, K)
\end{equation}
推出. 证明完成.
\end{proof}

\begin{example}[Klein--Gordon 方程]\label{eg:7.17}
回忆系统~\eqref{eq:7.17}--\eqref{eq:7.18}, 它借助~\eqref{eq:7.58} 被对角化. 定理~\ref{thm:7.16} 表明, 对光滑源项 $h$ 有
\begin{equation}\label{eq:7.86}
    \begin{aligned}
        \WF(U(t); 1) &= \WF\left( (SU(t))_1 \right) = \WF(\Lambda u(t) + D_t u(t)) \\
        &= \mathcal{C}_-(t, 0)\WF(\Lambda u_0 + u_1),
    \end{aligned}
\end{equation}
类似地
\begin{equation}\label{eq:7.87}
    \WF(U(t); 2) = \WF(\Lambda u(t) - D_t u(t)) = \mathcal{C}_+(t, 0)\WF(\Lambda u_0 - u_1);
\end{equation}
这里 $\mathcal{C}_\pm(t, 0)$ 是 $H_{\pm|\xi|} = \pm \frac{\xi}{|\xi|} \cdot \partial_x$ 的时间 $t$ 流. 作为~\eqref{eq:7.87} 的一个自洽性检验, 注意若 $(-D_t^2 + \Delta + 1)u = 0$, 则 $v \defeq \Lambda u - D_t u = -(D_t - \langle D \rangle)u$ 解方程 $(D_t + \langle D \rangle)v = 0$; 而 $v$ 的奇性沿 $\frac{\xi}{|\xi|} \cdot \partial_x$ 的传播已在例~\ref{eg:7.12} 中讨论过.
\end{example}

人们可以类似地研究波动方程~\eqref{eq:7.19} 解的波前集: 它们沿 $\partial_t + H_{\lambda_\pm(t)}$ 的积分曲线传播, 其中 $\lambda_\pm(t, x, \xi) \defeq \pm | \xi |_{g(t)_x^{-1}}$.



\section*{注记}
\addcontentsline{toc}{section}{注记}

我们对定理~\ref{thm:7.1}--\ref{thm:7.8} 的表述取自 Taylor~\cite[§§7.7--7.8]{Tay11b}. 证明定理~\ref{thm:7.1} 的另一种方法是对伴随问题 $(D_t + a(t, x, D_x)^*)v = h(t, x)$ 作先验能量估计, 再结合对偶论证. 这是 Hörmander~\cite[Chapter XXIII]{Hor07} 采取的处理方式, 在半经典情形下 Dyatlov--Zworski~\cite[Appendix E.5]{DZ19} 也是如此. 这些结果可以推广为关于 $m$ 阶严格双曲微分算子的半整体存在性与唯一性陈述; 从 $\mathcal{C}^0_t H^s_x$ 这类空间过渡到``时空''上的 Sobolev 空间, 需要用到具有独立的``时空''与``空间''正则性指标的中间空间. 我们在第~\ref{ch:09}~章讨论波动方程``时空''估计的一种更直接的途径.

伪微分算子在标量对称双曲方程之流下的共轭仍是伪微分算子 (定理~\ref{thm:7.6}), 这是 Egorov~\cite{Ego69} 的一个重要观察.

命题~\ref{prop:7.14} 中解耦论证的一个变体被 Taylor~\cite{Tay75} 用来证明带边界区域上双曲型方程的奇性反射.

\section*{习题}
\addcontentsline{toc}{section}{习题}

\begin{exer}[对称性假设的必要性]\label{ex:7.1}
在 $\R_t \times \R_x$ 上考虑方程 $(D_t + \ii D_x)u = 0$. (这是把 $u = u(t, x)$ 视为复变量 $t + \ii x$ 的函数时的 Cauchy--Riemann 方程.) 设 $u$ 是该方程的一个解, 且 $u \in \mathcal{C}^0(\R; H^s(\R^n))$. 关于 $u(0, x)$ 你能说什么? 由此断定定理~\ref{thm:7.1} 对 $a(t, x, D_x) = \ii D_x$ 不成立.
\end{exer}

\begin{exer}[可对称化的双曲系统]\label{ex:7.2}
设 $A \in \Smooth(\R_t; \Psi^1(\R^n)^{K \times K})$. 假定 $K \times K$ 系统
\begin{equation}\label{eq:7.88}
    \left\{
    \begin{aligned}
        (D_t + A(t)) u &= g \in \mathcal{C}^0\left( \R; H^s(\R^n; \CC^K) \right), \\
        u(0) &= u_0 \in H^s(\R^n; \CC^K)
    \end{aligned}
    \right.
\end{equation}
是\textbf{可对称化的双曲系统} (symmetrizable hyperbolic system): 意即存在一个 $K \times K$ 矩阵值符号 $s(t, x, \xi) \in S^0$, 它是正定的, 且使 $s(t, x, \xi)a(t, x, \xi)$ 模 $S^0$ 对称. 证明~\eqref{eq:7.88} 有唯一解
\begin{equation}\label{eq:7.89}
    u \in \mathcal{C}^0\left( \R; H^s(\R^n; \CC^K) \right) \cap \mathcal{C}^1\left( \R; H^{s-1}(\R^n; \CC^K) \right).
\end{equation}
\end{exer}

\begin{exer}[严格双曲系统, I]\label{ex:7.3}
我们称形如~\eqref{eq:7.88} 的 $K \times K$ 系统为\textbf{严格双曲的}, 若 $A$ 的主符号有一个齐次代表元 $a \in S^1_{\mathrm{hom}}$, 且它对一切 $t \in \R$, $x \in \R^n$ 与 $\xi \in \R^n \setminus \{0\}$ 有 $K$ 个互不相同的实特征值. (这与定义~\ref{def:7.13} 一致.) 证明严格双曲系统是可对称化的.
\end{exer}

\begin{exer}[严格双曲系统, II]\label{ex:7.4}
利用习题~\ref{ex:7.3} 证明定理~\ref{thm:7.16} 对一切严格双曲系统成立. (因此对称双曲性的假设并不必要.)
\end{exer}

\begin{exer}[高阶严格双曲方程]\label{ex:7.5}
在 $\R_t \times \R^n_x$ 上考虑 $m$ 阶算子
\begin{equation}\label{eq:7.90}
    L = D_t^m + \sum_{j=0}^{m-1} a_j(t, x, D_x) D_t^j,
\end{equation}
其中 $a_j(t, x, D_x) \in \operatorname{Diff}^{m-j}(\R^n_x)$ 光滑依赖于 $t$. 我们研究系统
\begin{equation}\label{eq:7.91}
    \left\{
    \begin{aligned}
        Lu &= g \in \mathcal{C}^0\left( \R_t; H^{s-m+1}(\R^n) \right), \\
        D_t^j u(0, \cdot) &= u_j \in H^{s-j}(\R^n), \qquad j = 0, \ldots, m-1.
    \end{aligned}
    \right.
\end{equation}
\begin{enumerate}
    \item 把~\eqref{eq:7.91} 化归为关于 $\CC^m$ 值函数 $(\Lambda^{m-1}u, \Lambda^{m-2}D_t u, \ldots, D_t^{m-1}u)$ 的一阶演化方程的 $m \times m$ 系统, 其中 $\Lambda = \langle D_x \rangle$.
    \item 给出 $L$ 的主符号上的一个条件, 使之等价于该 $m \times m$ 系统的严格双曲性.
    \item 在第 (2) 部分的严格双曲性假设下, 证明~\eqref{eq:7.91} 有唯一解 $u \in \bigcap_{j=0}^m \mathcal{C}^j(\R_t; H^{s-j}(\R^n))$.
\end{enumerate}
\end{exer}

